\documentclass[11pt,a4paper]{ctexart}
\usepackage[margin=2.2cm]{geometry}
\usepackage{xcolor}
\usepackage{enumitem}
\usepackage{amssymb}
\usepackage[normalem]{ulem}
\usepackage{titlesec}
\usepackage{hyperref}
\definecolor{navy}{HTML}{1A365D}
\definecolor{ink}{HTML}{1F2933}
\definecolor{muted}{HTML}{66788A}
\definecolor{soft}{HTML}{EEF4FA}
\definecolor{gold}{HTML}{8A5A12}
\hypersetup{colorlinks=true,linkcolor=navy,urlcolor=navy}
\setlist{nosep}
\titleformat{\section}{\Large\bfseries\color{navy}}{}{0pt}{}
\titleformat{\subsection}{\large\bfseries\color{navy}}{}{0pt}{}
\pagestyle{plain}

\title{\textbf{八上一般将来时与情态动词专项}\\[4pt]
\large 人教版（新课纲版）八年级上册 · 一般将来时 / 情态动词 can、could}
\author{初中英语资料站}
\date{2026-11-12}
\begin{document}
\maketitle

\section*{使用说明}
\begin{itemize}[leftmargin=1.4em]
\item 本讲义对应人教版（新课纲版）八年级上册语法要求：一般将来时与情态动词 can、could。
\item 共 4 讲语法、6 级 48 题。
\item 在线版：\url{https://1901151458-prog.github.io/english-materials/g8a-future/}
\end{itemize}

\section{一般将来时：will / shall}
一般将来时通常由「助动词 + 动词原形」构成，常与 tomorrow、next week、next year、in the future 等时间状语连用。
\begin{itemize}[leftmargin=1.4em]
\item We \underline{will come} to see you next week.
\item You \underline{will see} robots working everywhere in the future.
\item What \underline{will your life be like} in 20 years?
\end{itemize}
\textbf{缩写：}will → 'll；will not → won't。\\
\textbf{shall：}主语为 I / we 时可用于书面语。Shall we go to see the sci-fi film? / I shall write you a letter next month.

\section{一般将来时：be going to}
「be going to + 动词原形」表示将要发生的动作或安排，或打算、计划、决定要做的事。
\begin{itemize}[leftmargin=1.4em]
\item I'm \underline{going to make} a sandwich.
\item What \underline{are you going to do} tonight?
\item How \underline{are you going to improve} yourself?
\end{itemize}
\textbf{辨析：}有迹象表明即将发生（Look at the clouds! It is going to rain.）或已计划好的安排用 be going to；单纯预测或带意愿色彩用 will。

\section{情态动词 can}
\begin{center}
\begin{tabular}{p{3.0cm}p{10.4cm}}
\hline
\textbf{功能} & \textbf{例句}\\
\hline
能力 & \underline{Can} you speak English?\\
请求允许 & \underline{Can} I put some flowers on the table?\\
给予许可 & Yes, of course you \underline{can}.\\
提出请求 & \underline{Can} you help me hang up the photos?\\
提出建议 & We \underline{can} eat in a restaurant this evening.\\
\hline
\end{tabular}
\end{center}

\section{情态动词 could}
\begin{itemize}[leftmargin=1.4em]
\item \textbf{can 的过去式（能力）}：He \underline{could} sing very well when he was young.
\item \textbf{更委婉的请求}：\underline{Could} I use your mobile phone? / \underline{Could} you please sweep the floor?
\item \textbf{更委婉的建议}：We \underline{could} go fishing this weekend. / Maybe I \underline{could} give you a lift.
\end{itemize}
\textbf{要点：}表示请求时 could 比 can 更客气；但回答时通常仍用 can（Yes, you can.）。


\newpage
\section{分级练习（48 题）}
\noindent 按难度分级排列，先做基础，再做成分与扩展，最后做综合题。

\subsection*{第 1 级 · 一般将来时 · will}
\noindent\textcolor{gray}{will + 动词原形；won’t；’ll 缩写。}\par\smallskip
\begin{enumerate}[leftmargin=*,itemsep=4pt]
\item We \_\_\_ come to see you next week.
\begin{enumerate}[label=\Alph*.,leftmargin=2em,itemsep=1pt,topsep=2pt]
\item will
\item would
\item was
\item were
\end{enumerate}
\item You \_\_\_ see robots working everywhere in the future.
\begin{enumerate}[label=\Alph*.,leftmargin=2em,itemsep=1pt,topsep=2pt]
\item will
\item would
\item are
\item is
\end{enumerate}
\item What \_\_\_ your life be like in 20 years?
\begin{enumerate}[label=\Alph*.,leftmargin=2em,itemsep=1pt,topsep=2pt]
\item will
\item would
\item is
\item are
\end{enumerate}
\item Some jobs today \_\_\_ disappear in the coming years.
\begin{enumerate}[label=\Alph*.,leftmargin=2em,itemsep=1pt,topsep=2pt]
\item will
\item would
\item are
\item is
\end{enumerate}
\item I \_\_\_ tell you all about it.
\begin{enumerate}[label=\Alph*.,leftmargin=2em,itemsep=1pt,topsep=2pt]
\item will
\item would
\item am
\item is
\end{enumerate}
\item Tomorrow \_\_\_ be sunny and warm.
\begin{enumerate}[label=\Alph*.,leftmargin=2em,itemsep=1pt,topsep=2pt]
\item will
\item would
\item is
\item are
\end{enumerate}
\item He \_\_\_ arrive next Monday.
\begin{enumerate}[label=\Alph*.,leftmargin=2em,itemsep=1pt,topsep=2pt]
\item will
\item would
\item does
\item is
\end{enumerate}
\item They \_\_\_ not come tomorrow.
\begin{enumerate}[label=\Alph*.,leftmargin=2em,itemsep=1pt,topsep=2pt]
\item will
\item would
\item do
\item are
\end{enumerate}
\end{enumerate}

\subsection*{第 2 级 · 一般将来时 · be going to}
\noindent\textcolor{gray}{be going to + 动词原形，表示打算、计划或将要发生的事。}\par\smallskip
\begin{enumerate}[leftmargin=*,itemsep=4pt]
\item I’m \_\_\_ make a sandwich.
\begin{enumerate}[label=\Alph*.,leftmargin=2em,itemsep=1pt,topsep=2pt]
\item going to
\item going
\item go
\item gone
\end{enumerate}
\item What \_\_\_ you going to do tonight?
\begin{enumerate}[label=\Alph*.,leftmargin=2em,itemsep=1pt,topsep=2pt]
\item are
\item is
\item do
\item will
\end{enumerate}
\item She \_\_\_ going to read more books about AI.
\begin{enumerate}[label=\Alph*.,leftmargin=2em,itemsep=1pt,topsep=2pt]
\item is
\item are
\item do
\item will
\end{enumerate}
\item We \_\_\_ going to visit the museum.
\begin{enumerate}[label=\Alph*.,leftmargin=2em,itemsep=1pt,topsep=2pt]
\item are
\item is
\item do
\item will
\end{enumerate}
\item How \_\_\_ you going to improve yourself?
\begin{enumerate}[label=\Alph*.,leftmargin=2em,itemsep=1pt,topsep=2pt]
\item are
\item is
\item do
\item will
\end{enumerate}
\item He \_\_\_ going to buy a new bike.
\begin{enumerate}[label=\Alph*.,leftmargin=2em,itemsep=1pt,topsep=2pt]
\item is
\item are
\item does
\item will
\end{enumerate}
\item It \_\_\_ going to rain soon.
\begin{enumerate}[label=\Alph*.,leftmargin=2em,itemsep=1pt,topsep=2pt]
\item is
\item are
\item does
\item will
\end{enumerate}
\item They \_\_\_ going to have a party.
\begin{enumerate}[label=\Alph*.,leftmargin=2em,itemsep=1pt,topsep=2pt]
\item are
\item is
\item do
\item will
\end{enumerate}
\end{enumerate}

\subsection*{第 3 级 · 一般将来时 · 综合}
\noindent\textcolor{gray}{will / shall / be going to 的辨析与时间状语。}\par\smallskip
\begin{enumerate}[leftmargin=*,itemsep=4pt]
\item Shall we \_\_\_ to see the sci-fi film?
\begin{enumerate}[label=\Alph*.,leftmargin=2em,itemsep=1pt,topsep=2pt]
\item go
\item going
\item went
\item gone
\end{enumerate}
\item I shall \_\_\_ you a letter next month.
\begin{enumerate}[label=\Alph*.,leftmargin=2em,itemsep=1pt,topsep=2pt]
\item write
\item writing
\item wrote
\item written
\end{enumerate}
\item will not 的缩写形式是 \_\_\_。
\begin{enumerate}[label=\Alph*.,leftmargin=2em,itemsep=1pt,topsep=2pt]
\item won’t
\item willn’t
\item wonot
\item wannot
\end{enumerate}
\item will 在口语中的缩写形式是 \_\_\_。
\begin{enumerate}[label=\Alph*.,leftmargin=2em,itemsep=1pt,topsep=2pt]
\item ’ll
\item ’d
\item ’s
\item ’re
\end{enumerate}
\item My parents \_\_\_ buy me a computer next week.
\begin{enumerate}[label=\Alph*.,leftmargin=2em,itemsep=1pt,topsep=2pt]
\item will
\item would
\item do
\item are
\end{enumerate}
\item Look at the dark clouds! It \_\_\_ rain.
\begin{enumerate}[label=\Alph*.,leftmargin=2em,itemsep=1pt,topsep=2pt]
\item is going to
\item will
\item would
\item does
\end{enumerate}
\item We \_\_\_ have a picnic this Sunday. We’ve planned it.
\begin{enumerate}[label=\Alph*.,leftmargin=2em,itemsep=1pt,topsep=2pt]
\item are going to
\item will
\item would
\item do
\end{enumerate}
\item What \_\_\_ you do tomorrow?
\begin{enumerate}[label=\Alph*.,leftmargin=2em,itemsep=1pt,topsep=2pt]
\item will
\item would
\item do
\item are
\end{enumerate}
\end{enumerate}

\subsection*{第 4 级 · 情态动词 can}
\noindent\textcolor{gray}{能力、请求允许、给予许可、提出请求和建议。}\par\smallskip
\begin{enumerate}[leftmargin=*,itemsep=4pt]
\item \_\_\_ you speak English?
\begin{enumerate}[label=\Alph*.,leftmargin=2em,itemsep=1pt,topsep=2pt]
\item Can
\item Should
\item Must
\item Need
\end{enumerate}
\item Can you help me \_\_\_ up the photos?
\begin{enumerate}[label=\Alph*.,leftmargin=2em,itemsep=1pt,topsep=2pt]
\item hang
\item hanging
\item hung
\item hangs
\end{enumerate}
\item A: Can I put some flowers on the table? B: Yes, of course you \_\_\_.
\begin{enumerate}[label=\Alph*.,leftmargin=2em,itemsep=1pt,topsep=2pt]
\item can
\item could
\item must
\item need
\end{enumerate}
\item We \_\_\_ eat in a restaurant this evening.
\begin{enumerate}[label=\Alph*.,leftmargin=2em,itemsep=1pt,topsep=2pt]
\item can
\item must
\item need
\item should
\end{enumerate}
\item Can you \_\_\_ me with this box?
\begin{enumerate}[label=\Alph*.,leftmargin=2em,itemsep=1pt,topsep=2pt]
\item help
\item helping
\item helped
\item helps
\end{enumerate}
\item \_\_\_ I use your pen?
\begin{enumerate}[label=\Alph*.,leftmargin=2em,itemsep=1pt,topsep=2pt]
\item Can
\item Must
\item Need
\item Should
\end{enumerate}
\item She can \_\_\_ very well.
\begin{enumerate}[label=\Alph*.,leftmargin=2em,itemsep=1pt,topsep=2pt]
\item sing
\item singing
\item sang
\item sings
\end{enumerate}
\item You can \_\_\_ here if you like.
\begin{enumerate}[label=\Alph*.,leftmargin=2em,itemsep=1pt,topsep=2pt]
\item sit
\item sitting
\item sat
\item sits
\end{enumerate}
\end{enumerate}

\subsection*{第 5 级 · 情态动词 could}
\noindent\textcolor{gray}{can 的过去式（能力）；更委婉的请求与建议。}\par\smallskip
\begin{enumerate}[leftmargin=*,itemsep=4pt]
\item He \_\_\_ sing very well when he was young.
\begin{enumerate}[label=\Alph*.,leftmargin=2em,itemsep=1pt,topsep=2pt]
\item could
\item can
\item must
\item need
\end{enumerate}
\item \_\_\_ you please sweep the floor?
\begin{enumerate}[label=\Alph*.,leftmargin=2em,itemsep=1pt,topsep=2pt]
\item Could
\item Must
\item Need
\item Should
\end{enumerate}
\item We \_\_\_ go fishing this weekend.
\begin{enumerate}[label=\Alph*.,leftmargin=2em,itemsep=1pt,topsep=2pt]
\item could
\item must
\item need
\item should
\end{enumerate}
\item Maybe I \_\_\_ give you a lift to the supermarket.
\begin{enumerate}[label=\Alph*.,leftmargin=2em,itemsep=1pt,topsep=2pt]
\item could
\item must
\item need
\item should
\end{enumerate}
\item I \_\_\_ ride a bike when I was six.
\begin{enumerate}[label=\Alph*.,leftmargin=2em,itemsep=1pt,topsep=2pt]
\item could
\item can
\item must
\item need
\end{enumerate}
\item \_\_\_ you help me with this heavy bag?
\begin{enumerate}[label=\Alph*.,leftmargin=2em,itemsep=1pt,topsep=2pt]
\item Could
\item Must
\item Need
\item Should
\end{enumerate}
\item She \_\_\_ play the piano when she was young.
\begin{enumerate}[label=\Alph*.,leftmargin=2em,itemsep=1pt,topsep=2pt]
\item could
\item can
\item must
\item need
\end{enumerate}
\item You \_\_\_ take a taxi if you’re in a hurry.
\begin{enumerate}[label=\Alph*.,leftmargin=2em,itemsep=1pt,topsep=2pt]
\item could
\item must
\item need
\item should
\end{enumerate}
\end{enumerate}

\subsection*{第 6 级 · 综合提升}
\noindent\textcolor{gray}{一般将来时与 can / could 混合辨析。}\par\smallskip
\begin{enumerate}[leftmargin=*,itemsep=4pt]
\item We \_\_\_ go to the park tomorrow if it’s fine.
\begin{enumerate}[label=\Alph*.,leftmargin=2em,itemsep=1pt,topsep=2pt]
\item will
\item would
\item are
\item is
\end{enumerate}
\item He \_\_\_ play basketball when he was 10.
\begin{enumerate}[label=\Alph*.,leftmargin=2em,itemsep=1pt,topsep=2pt]
\item could
\item can
\item must
\item need
\end{enumerate}
\item \_\_\_ we go to the zoo this weekend?
\begin{enumerate}[label=\Alph*.,leftmargin=2em,itemsep=1pt,topsep=2pt]
\item Shall
\item Should
\item Must
\item Need
\end{enumerate}
\item I’m sure he \_\_\_ come tomorrow.
\begin{enumerate}[label=\Alph*.,leftmargin=2em,itemsep=1pt,topsep=2pt]
\item will
\item would
\item does
\item is
\end{enumerate}
\item \_\_\_ you speak Japanese?
\begin{enumerate}[label=\Alph*.,leftmargin=2em,itemsep=1pt,topsep=2pt]
\item Can
\item Must
\item Need
\item Should
\end{enumerate}
\item The train \_\_\_ leave at 8:00 tomorrow morning.
\begin{enumerate}[label=\Alph*.,leftmargin=2em,itemsep=1pt,topsep=2pt]
\item will
\item would
\item does
\item is
\end{enumerate}
\item We \_\_\_ going to have a test next week.
\begin{enumerate}[label=\Alph*.,leftmargin=2em,itemsep=1pt,topsep=2pt]
\item are
\item is
\item do
\item will
\end{enumerate}
\item It \_\_\_ going to snow tonight.
\begin{enumerate}[label=\Alph*.,leftmargin=2em,itemsep=1pt,topsep=2pt]
\item is
\item are
\item does
\item will
\end{enumerate}
\end{enumerate}

\newpage
\section*{答案与解析}
\small
\begin{enumerate}[leftmargin=1.6em,itemsep=2pt]
\item \textbf{F1-1 A}：next week 表将来，用 will；教材例句。
\item \textbf{F1-2 A}：in the future 表将来，用 will；教材例句。
\item \textbf{F1-3 A}：in 20 years 表将来；教材例句。
\item \textbf{F1-4 A}：in the coming years 表将来；教材例句。
\item \textbf{F1-5 A}：表示带意愿色彩的将来常用 will；教材例句。
\item \textbf{F1-6 A}：tomorrow 表将来，用 will + 动词原形。
\item \textbf{F1-7 A}：next Monday 表将来，用 will。
\item \textbf{F1-8 A}：一般将来时否定：will not = won’t。
\item \textbf{F2-1 A}：be going to + 动词原形；教材例句。
\item \textbf{F2-2 A}：主语 you 用 are going to；教材例句。
\item \textbf{F2-3 A}：主语 she 用 is going to。
\item \textbf{F2-4 A}：主语 we 用 are going to。
\item \textbf{F2-5 A}：主语 you 用 are；教材例句。
\item \textbf{F2-6 A}：主语 he 用 is going to。
\item \textbf{F2-7 A}：主语 it 用 is going to。
\item \textbf{F2-8 A}：主语 they 用 are going to。
\item \textbf{F3-1 A}：shall + 动词原形；教材例句。
\item \textbf{F3-2 A}：shall + 动词原形；教材例句。
\item \textbf{F3-3 A}：will not 常缩略为 won’t。
\item \textbf{F3-4 A}：will 常缩略为 ’ll。
\item \textbf{F3-5 A}：next week 表将来，用 will。
\item \textbf{F3-6 A}：有迹象表明将要发生，用 be going to。
\item \textbf{F3-7 A}：已经计划好的安排用 be going to。
\item \textbf{F3-8 A}：tomorrow 表将来，用 will + 动词原形。
\item \textbf{F4-1 A}：can 表示能力；教材例句。
\item \textbf{F4-2 A}：情态动词后接动词原形；教材例句。
\item \textbf{F4-3 A}：给予许可用 can；教材例句。
\item \textbf{F4-4 A}：can 提出建议；教材例句。
\item \textbf{F4-5 A}：情态动词后接动词原形。
\item \textbf{F4-6 A}：Can I ...? 请求允许。
\item \textbf{F4-7 A}：情态动词后接动词原形。
\item \textbf{F4-8 A}：情态动词后接动词原形。
\item \textbf{F5-1 A}：could 表示过去的能力；教材例句。
\item \textbf{F5-2 A}：Could you please ...? 更委婉的请求；教材例句。
\item \textbf{F5-3 A}：could 提出建议；教材例句。
\item \textbf{F5-4 A}：could 提出建议，语气委婉；教材例句。
\item \textbf{F5-5 A}：过去的能力用 could。
\item \textbf{F5-6 A}：Could you ...? 委婉请求。
\item \textbf{F5-7 A}：过去的能力用 could。
\item \textbf{F5-8 A}：could 提出建议。
\item \textbf{F6-1 A}：tomorrow 表将来，用 will。
\item \textbf{F6-2 A}：过去的能力用 could。
\item \textbf{F6-3 A}：Shall we ...? 提出建议，主语为 we 时用 shall。
\item \textbf{F6-4 A}：tomorrow 表将来，用 will。
\item \textbf{F6-5 A}：can 表示能力。
\item \textbf{F6-6 A}：tomorrow morning 表将来，用 will。
\item \textbf{F6-7 A}：主语 we 用 are going to。
\item \textbf{F6-8 A}：主语 it 用 is going to。
\end{enumerate}

\end{document}
